% ------------------------------------------------------------
% CHAPTER 36
% ------------------------------------------------------------

\chapter{Active Control, Agency, and Adaptive Measurement}

The delta–sigma modulator is often described as a converter, a device whose only purpose is to translate an analog signal into a discrete sequence. Yet the underlying structure of the modulator suggests a far more general interpretation. A delta–sigma loop does not merely receive a signal; it actively regulates its internal state to minimize error and preserve coherence. It predicts, corrects, and stabilizes. It performs a kind of work to maintain the alignment between its internal dynamics and the external field it measures. These properties are not incidental. They are characteristic of systems that possess a minimal form of agency. In this chapter, we develop the formal equivalence between delta–sigma modulation and active control. We show that when viewed through the lenses of entropy, geometry, and information flow, the modulator becomes a prototype of an adaptive measurement agent that continually acts to shape the uncertainty it encounters.

Agency may be defined as the capacity of a system to maintain certain internal invariants despite external disturbances. For biological and cognitive systems, these invariants correspond to viability conditions, homeostatic bounds, or semantic coherence. For engineered control systems, they correspond to stability criteria or tracking performance. For the delta–sigma modulator, the invariant is the relationship between the input signal and the internal state: the integrator must remain bounded, the feedback must remain coherent, and the shaped error spectrum must remain within predictable limits. If any of these invariants are lost, the modulator ceases to function. Its entire architecture is designed to resist the accumulation of entropy in sensitive degrees of freedom. This resistance is the hallmark of agency in the most minimal sense.

One can formalize this claim by examining the modulator’s internal state dynamics. The internal state evolves as
\[
u[n+1] = u[n] + x[n] - y[n],
\]
which expresses the difference between prediction and outcome. The modulator’s sole objective is to keep the error \(x[n] - y[n]\) small enough that the internal state does not diverge. The feedback law embodied in the loop filter chooses \(y[n]\) so as to counteract the drift induced by the input. Importantly, the modulator chooses its output symbol based not only on the current input but on the entire internal history captured in \(u[n]\). This dependence encodes a memory of prior mismatches, allowing the modulator to act with respect to a long-term objective: preservation of internal coherence. In this sense, the modulator is a control agent whose actuator is the quantizer and whose policy is the feedback law implemented by the loop filter.

Entropy considerations provide further justification for interpreting the modulator as an agent. The quantization step injects entropy into the system by collapsing a continuous set of possible states onto a discrete symbol. The loop filter serves as the mechanism by which the modulator routes this entropy into regions of representational space where it does not threaten coherence. Without such routing, the error would accumulate uncontrollably, saturating the integrator and causing instability. The modulator therefore must act to remove the entropy that measurement injects. This dynamic mirrors the thermodynamic role played by biological and cognitive agents, which must counteract entropy production to maintain stable function. Agency emerges from the interplay between measurement-induced entropy and active regulation.

When the modulator is generalized to higher order, the analogy becomes even more striking. The internal state becomes a high-dimensional field \(U(n)\) living in a manifold whose geometry is determined by the filter coefficients. The evolution equation
\[
U[n+1] = A U[n] + B x[n] - C y[n]
\]
encodes a drift field that reflects the system’s priors about how the input is expected to evolve. The modulator must then choose \(y[n]\) so that the internal trajectory remains confined to the stable region of this manifold. This is precisely the structure of a controlled dynamical system performing active inference. The internal state represents beliefs or predictions; the quantized output represents actions taken to correct discrepancies; and the feedback represents the coupling between belief and action that maintains coherence.

To express this structure in the language of active control, one may define an internal free-energy functional
\[
F(U[n]) = \frac{1}{2} \left\| A U[n] - B x[n] \right\|^2 + \lambda \, S(U[n]),
\]
where the first term penalizes mismatch between predicted and observed inputs and the second term penalizes the growth of entropy in the internal reservoir. The modulation loop seeks to minimize \(F\). The discrete symbol \(y[n]\) is chosen to reduce the free-energy gradient, and the residual error produced by this choice becomes the quantization noise that must be routed away from critical degrees of freedom. The modulator’s behavior therefore corresponds to a gradient-descent policy on an internal energy landscape. The quantizer’s discontinuity simply forces this policy to operate in a space where actions are discrete, but the optimization principle is continuous.

From this viewpoint, delta–sigma modulation becomes a specific instance of a much broader class of systems that enact minimal agency. A thermostat, a servo controller, and a simple neuron all operate according to the same principle: observe a field, collapse the observation into a discrete state, integrate error, and take actions that maintain invariants. The modulator is distinguished by its use of oversampling, which gives it access to a high temporal resolution, allowing it to respond rapidly to fluctuations and distribute entropy finely across time. This oversampling is the temporal analogue of the richness of a sensory receptor array. It affords the modulator the ability to enact its control policy in a fine-grained manner.

Adaptive measurement arises naturally from this adaptive control interpretation. A modulator can adjust its loop filter coefficients, quantizer thresholds, or dither characteristics to optimize performance under changing input statistics. Such adaptation is mathematically equivalent to adjusting the policy parameters in an active control agent. The modulator, when equipped with adaptive mechanisms, becomes capable of learning how to distribute entropy more effectively. It can change the geometry of its internal state space so that prediction error flows into more favorable regions. It can alter its collapse thresholds to better match the distribution of its inputs. These changes reflect the emergence of a rudimentary form of plasticity analogous to synaptic plasticity in neurons.

One may go further and observe that a delta–sigma modulator equipped with adaptive parameters behaves as a minimal active inference agent. It maintains an internal generative model of its inputs through the dynamics of its integrators. It predicts the incoming signal based on past dynamics. It collapses its representation into discrete symbols, which serve as both output and internal corrective action. It routes entropy into harmless regions of the representational space through noise shaping. And it acts to maintain stability by minimizing internal free energy. Although the modulator’s symbolic outputs are simple, its internal regulation reflects the essential operations of agency: prediction, correction, coherence maintenance, and entropy control.

This viewpoint unifies engineered control systems with biological and cognitive agents. It suggests that agency is not a mysterious emergent property but a natural consequence of systems organized around measurement, collapse, entropy routing, and feedback. The delta–sigma modulator thus becomes a clean mathematical model of the minimal structure required for agency. It operates at the boundary between continuous fields and discrete symbols, transforms entropy into structure, and maintains its own coherence despite continual perturbations.

The next chapter builds on this foundation by exploring how delta–sigma dynamics, RSVP theory, and recursive inference combine to produce coherent systems that operate across scales. It shows how the principles developed here generalize to complex networks that maintain shared identity, compressed representation, and stable semantics through a hierarchical geometry of feedback loops.

